\subsection{最初に必要な用語}
\par\smallskip\noindent\refstepcounter{table}\label{tab:ch13-01}表 \thetable: 最初に必要な用語の整理\par\smallskip
表 \thetable は、最初に必要な用語の整理を比較・確認しやすい形で整理したものである。\par\smallskip
\begin{longtable}{>{\raggedright\arraybackslash}p{0.25\linewidth}>{\raggedright\arraybackslash}p{0.7\linewidth}}
\toprule
用語 & 初学者向けの説明 \\
\midrule
\endhead
ソフトウェアセキュリティ & 情報やデータを、権限に応じて適切に守るソフトウェアの品質特性である。 \\
情報セキュリティ & 情報の機密性、完全性、可用性を保つことを中心にした考え方である。 \\
サイバーセキュリティ & サイバー空間におけるリスクを許容できる水準に保つため、人、組織、社会、国家を守る活動である。 \\
\term{脆弱性}{Vulnerability} & 攻撃者が悪用できる誤り、弱点、設定不備である。 \\
攻撃者 & システムを不正利用しようとする人、組織、自動化されたプログラムなどである。 \\
\term{脅威}{Threat} & 攻撃や事故によって、守るべき資産に損害を与える可能性である。 \\
資産 & 守る価値がある情報、機能、サービス、鍵、ログ、信用、設備などである。 \\
\term{SDLC}{Software Development Life Cycle} & セキュア開発ライフサイクル。開発の各段階にセキュリティ活動を組み込む考え方である。 \\
DevSecOps & 開発、セキュリティ、運用を統合し、継続的な自動化と改善の中でセキュリティを扱う考え方である。 \\
静的分析 & プログラムを実行せず、ソースコードやバイナリを調べて弱点を探す方法である。 \\
\term{動的テスト}{Dynamic Testing} & ソフトウェアを実行し、入力や操作に対する振る舞いから弱点を探す方法である。 \\
ペネトレーションテスト & 許可された範囲で実際の攻撃に近い操作を行い、侵入可能性や悪用経路を確認する試験である。 \\
\bottomrule
\end{longtable}
